\section{A2C Algorithm Applied and Analysis}

For each of the three shares, an Advantage Actor-Critic (A2C) agent was trained and
evaluated against a simple buy-and-hold strategy. Per share, five independently trained
agents were compared with the market. The tables below report the average returns of the
trained agents, the average returns of the market (buy and hold), and the average
percentage of cases in which the agent performed better than the market.

\subsection{Boeing}

\begin{table}[H]
  \centering
  \caption{Boeing -- average returns of agent vs.\ market.}
  \label{tab:boeing}
  \begin{tabular}{lrrrrr}
    \toprule
     & Agent 1 & Agent 2 & Agent 3 & Agent 4 & Agent 5\\
    \midrule
    Trained agent   & 1.151250 & 0.982355 & 1.268870 & 1.211902 & 1.220214\\
    Market (B\&H)   & 1.220214 & 1.194680 & 1.411054 & 1.325877 & 1.139906\\
    Agent better (\%) & 0.55 & 0.10 & 0.40 & 0.80 & 0.40\\
    \bottomrule
  \end{tabular}
\end{table}

\begin{figure}[H]
  \centering
  \includegraphics[width=0.78\textwidth]{z043_1.png}
  \caption{Boeing -- cumulative return ratio: buy and hold vs.\ trained agent.}
\end{figure}

\begin{figure}[H]
  \centering
  \includegraphics[width=0.78\textwidth]{z044_1.png}
  \caption{Boeing -- actions chosen by the trained agents over time.}
\end{figure}

\paragraph{Commentary on the result for Boeing.}
Given those numbers, it can be stated that the A2C algorithm applied to the Boeing share
would not result in an overall better performance than the simple buy-and-hold strategy.
Only in 45\% of the cases (on average) would the trained agent show a better performance.
This is mainly due to one entrained agent which leads to better results than the market in
only 10\% of the cases. The analysis of the performed actions shows significantly
different values, especially for the values for \emph{short} and \emph{cash}, distributed
among the trained agents. It can therefore be assumed that the trained agent applied to
this stock leads to inconsistent decisions and therefore to inconsistent and rather less
promising results.

\subsection{Airbus}

\begin{table}[H]
  \centering
  \caption{Airbus -- average returns of agent vs.\ market.}
  \label{tab:airbus}
  \begin{tabular}{lrrrrr}
    \toprule
     & Agent 1 & Agent 2 & Agent 3 & Agent 4 & Agent 5\\
    \midrule
    Trained agent   & 1.430267 & 1.570172 & 2.007212 & 1.276262 & 1.038152\\
    Market (B\&H)   & 1.217388 & 1.188968 & 1.168107 & 1.195885 & 1.237382\\
    Agent better (\%) & 0.75 & 0.60 & 0.80 & 0.60 & 0.20\\
    \bottomrule
  \end{tabular}
\end{table}

\begin{figure}[H]
  \centering
  \includegraphics[width=0.78\textwidth]{z058_1.png}
  \caption{Airbus -- cumulative return ratio: buy and hold vs.\ trained agent.}
\end{figure}

\begin{figure}[H]
  \centering
  \includegraphics[width=0.78\textwidth]{z059_1.png}
  \caption{Airbus -- actions chosen by the trained agents over time.}
\end{figure}

\paragraph{Commentary on the result for Airbus.}
Given those numbers, it can be stated that the A2C algorithm applied to the Airbus share
would result in an overall better performance than the simple buy-and-hold strategy. In
60\% of the cases (on average) the trained agent would show a better performance. The
analysis of the performed actions shows similar values for four trained agents; only one
trained agent differs a lot in terms of \emph{cash}. It can therefore be assumed that the
trained agent applied to this share could lead to consistent and promising results.

\subsection{Toyota}

\begin{table}[H]
  \centering
  \caption{Toyota -- average returns of agent vs.\ market.}
  \label{tab:toyota}
  \begin{tabular}{lrrrrr}
    \toprule
     & Agent 1 & Agent 2 & Agent 3 & Agent 4 & Agent 5\\
    \midrule
    Trained agent   & 1.160707 & 1.442139 & 1.185587 & 1.263972 & 1.272148\\
    Market (B\&H)   & 0.994851 & 1.014057 & 1.044513 & 1.070050 & 1.003871\\
    Agent better (\%) & 0.80 & 0.70 & 0.80 & 0.85 & 0.75\\
    \bottomrule
  \end{tabular}
\end{table}

\begin{figure}[H]
  \centering
  \includegraphics[width=0.78\textwidth]{z073_1.png}
  \caption{Toyota -- cumulative return ratio: buy and hold vs.\ trained agent.}
\end{figure}

\begin{figure}[H]
  \centering
  \includegraphics[width=0.78\textwidth]{z074_1.png}
  \caption{Toyota -- actions chosen by the trained agents over time.}
\end{figure}

\paragraph{Commentary on the result for Toyota.}
Given those numbers, it can be stated that the A2C algorithm applied to the Toyota share
would result in an overall better performance than the simple buy-and-hold strategy. In
78\% of the cases (on average) the trained agent would show a better performance. The
analysis of the performed actions shows different values for all trained agents. It can
therefore be assumed that the trained agent applied to this share could lead to
inconsistent behavior but could also lead to overall promising results.
